\subsection{道路连通空间}

定义 3. --- 若对拓扑空间 X 的每一对点 \((x,y)\)，都存在以 x 为起点、y 为终点的道路，则称 X 是道路连通的。

例 1. --- R 的每个区间、数值空间的每个凸子集，或更一般地，R 上拓扑向量空间的每个凸子集，都是道路连通的。

命题 3. --- 道路连通空间在连续映射下的像是道路连通的。

设 X 为道路连通空间，\(f: X \to Y\) 为连续映射。设 x 与 y 为 \(f(X)\) 的两点；设 \(x'\) 与 \(y'\) 为 X 中满足 \(f(x') = x\) 与 \(f(y') = y\) 的点。存在 X 中以 \(x'\) 为起点、\(y'\) 为终点的道路 c。于是 \(f \circ c\) 是 \(f(X)\) 中以 x 为起点、y 为终点的道路。

命题 4. --- 道路连通拓扑空间是连通的。

由于 R 的每个区间都连通，道路的像是连通的（TG, I, p. 82, 命题 4）。空空间是连通的。由于非空道路连通空间是从其中一点出发的道路的像的并，它是连通的（TG, I, p. 81, 命题 2）。

注意到（III, p. 256, 注 1）X 中的道路与从缩化为一点的拓扑空间 P 到 X 中的映射之间的同伦的等同，III, p. 230 的命题 1 蕴涵：

命题 5. --- 在拓扑空间中，关系“存在以 x 为起点、y 为终点的道路”是一个等价关系。

拓扑空间 X 的道路连通分支指上述关系的等价类。道路连通分支是道路连通空间。设 x 为 X 的一点。x 的道路连通分支是 X 中包含 x 的道路连通子空间的并。它也是 X 中以 x 为起点的道路的像的并。

例。--- 2) 交非空的一族道路连通集的并是道路连通的（参见 TG, I, p. 81, 命题 2）。

\begin{enumerate}
\def\labelenumi{\arabic{enumi})}
\setcounter{enumi}{2}
\tightlist
\item
  数值空间（或更一般地，R 上的拓扑向量空间）的在其中一点处为星形的子集是道路连通的。
\end{enumerate}

设 X 为拓扑空间。用 \(\pi_{0}(X)\) 表示 X 的道路连通分支组成的集合。设 P 为缩化为一点的拓扑空间。从 X 到 {[}P; X{]} 的映射把 X 的每一点 x 对应到同伦类 \(\varphi_{x}\)，即像为 x 的映射 \(f_{x}: P \to X\) 的同伦类，它通过取商定义了一个从 \(\pi_{0}(X)\) 到 {[}P; X{]} 上的双射，称为典范双射。我们有 \(\pi_{0}(\varnothing) = \varnothing\)。为使非空拓扑空间道路连通，必须而且只需 \(\pi_{0}(X)\) 是单元素集。

设 X 与 Y 为拓扑空间，\(f\colon\mathrm{X}\to\mathrm{Y}\) 为连续映射。X 的每个道路连通分支 C 的像 \(f(\mathrm{C})\) 是道路连通的（III, p. 258, 命题 3），因此包含于 Y 的某个道路连通分支。用 \(\pi_{0}(f)\colon\pi_{0}(\mathrm{X})\to\pi_{0}(\mathrm{Y})\) 表示这样的映射：它把 X 的道路连通分支 C 对应到满足 \(f(\mathrm{C})\subset\mathrm{C}'\) 的唯一的 Y 的道路连通分支 C'。若把 \(\pi_0(\mathrm{X})\) 与 \(\pi_0(\mathrm{Y})\) 分别等同于 {[}P; X{]} 与 {[}P; Y{]}，则 \(\pi_0(f)\) 等同于映射 \(\chi \mapsto [f] \circ \chi\)（从 {[}P; X{]} 到 {[}P; Y{]}）。特别地，若 \(f\) 与 \(g\) 是从 X 到 Y 的互相同伦的映射，则有 \(\pi_0(f) = \pi_0(g)\)（III, p. 230, 命题 2）。

设 Z 为拓扑空间，\(g: Y \to Z\) 为连续映射；我们有 \(\pi_{0}(g \circ f) = \pi_{0}(g) \circ \pi_{0}(f)\)。特别地，若 Z = X 且 f 与 g 互为同伦逆，则有 \(\pi_{0}(g) \circ \pi_{0}(f) = \pi_{0}(\mathrm{Id}_{\mathrm{X}}) = \mathrm{Id}_{\pi_{0}(\mathrm{X})}\) 与 \(\pi_{0}(f) \circ \pi_{0}(g) = \pi_{0}(\mathrm{Id}_{\mathrm{Y}}) = \mathrm{Id}_{\pi_{0}(\mathrm{Y})}\)，这证明 \(\pi_{0}(f)\) 与 \(\pi_{0}(g)\) 是互逆的双射。因此，与道路连通空间同伦等价的空间本身是道路连通的。特别地，与一点同伦等价的空间是道路连通的。

命题 6. --- 设 \((\mathrm{Y}_{j})_{j\in\mathrm{J}}\) 为一族拓扑空间。映射

\[
(\pi_ {0} (\mathrm{pr} _ {j})) \colon \pi_ {0} (\prod_ {j \in \mathrm{J}} \mathrm{Y} _ {j}) \to \prod_ {j \in \mathrm{J}} \pi_ {0} (\mathrm{Y} _ {j})
\]

是双射。特别地，一族道路连通空间的积空间是道路连通的。

这由 III, p. 231 的命题 3 得出，其中取 X 为缩化为单独一点的空间。

